% Figure: analytic M/M/c prediction against a discrete-event Monte-Carlo
% simulation, with the staffing-cost optimisation curve alongside.
\begin{figure}[htbp]
\centering
\begin{tikzpicture}[font=\small]
% ===== Left: analytic vs simulation, mean wait vs number of servers =====
\begin{scope}[xshift=0cm]
\draw[->, thick, draw=engblack!70] (0,0) -- (6.0,0)
  node[right, font=\scriptsize] {servers $c$};
\draw[->, thick, draw=engblack!70] (0,0) -- (0,4.0)
  node[above, font=\scriptsize] {mean wait $W_q$ (min)};
\foreach \xt/\xv in {1/3, 2/4, 3/5, 4/6, 5/7} {
  \draw[draw=engblack!70] (\xt,0) -- (\xt,-0.07);
  \node[font=\scriptsize, anchor=north] at (\xt,-0.09) {\xv};
}
\foreach \yt/\yv in {0/0, 1/10, 2/20, 3/30} {
  \draw[draw=engblack!70] (0,\yt) -- (-0.07,\yt);
  \node[font=\scriptsize, anchor=east] at (-0.09,\yt) {\yv};
}
% Analytic Erlang-C wait (decreasing, convex): tabulated points as a curve
\draw[very thick, draw=engblue!80] plot[smooth]
  coordinates {(1,3.6) (2,1.45) (3,0.62) (4,0.27) (5,0.11)};
% Simulation points with error bars (close to analytic)
\foreach \x/\y/\e in {1/3.45/0.30, 2/1.52/0.18, 3/0.58/0.10, 4/0.30/0.06, 5/0.10/0.04} {
  \draw[draw=engred!85, thick] (\x,{\y-\e}) -- (\x,{\y+\e});
  \fill[engred!85] (\x,\y) circle (1.6pt);
}
\node[font=\small, anchor=south west] at (0.1, 3.7)
  {\textbf{(a) prediction vs simulation}};
\draw[very thick, draw=engblue!80] (3.6,3.2) -- (4.0,3.2);
\node[anchor=west, font=\scriptsize] at (4.05,3.2) {Erlang C};
\fill[engred!85] (3.8,2.85) circle (1.6pt);
\node[anchor=west, font=\scriptsize] at (4.05,2.85) {sim.\ $\pm$SE};
\end{scope}
% ===== Right: total cost vs number of servers (the optimisation) =====
\begin{scope}[xshift=8.2cm]
\draw[->, thick, draw=engblack!70] (0,0) -- (6.0,0)
  node[right, font=\scriptsize] {servers $c$};
\draw[->, thick, draw=engblack!70] (0,0) -- (0,4.0)
  node[above, font=\scriptsize] {cost per hour (scaled)};
\foreach \xt/\xv in {1/3, 2/4, 3/5, 4/6, 5/7} {
  \draw[draw=engblack!70] (\xt,0) -- (\xt,-0.07);
  \node[font=\scriptsize, anchor=north] at (\xt,-0.09) {\xv};
}
% staffing cost: linear rising
\draw[thick, draw=enggray, dashed] plot coordinates
  {(1,0.8) (2,1.2) (3,1.6) (4,2.0) (5,2.4)};
\node[anchor=west, font=\scriptsize, text=enggray] at (4.6,2.5) {staff};
% waiting cost: convex falling (proportional to W_q * lambda)
\draw[thick, draw=engamber, dash dot] plot[smooth] coordinates
  {(1,3.6) (2,1.7) (3,0.85) (4,0.45) (5,0.25)};
\node[anchor=west, font=\scriptsize, text=engamber] at (4.7,0.35) {wait};
% total cost: U-shape with minimum near c=4
\draw[very thick, draw=engblue!85] plot[smooth] coordinates
  {(1,4.4) (2,2.9) (3,2.45) (4,2.45) (5,2.65)};
\node[anchor=west, font=\scriptsize, text=engblue!85] at (1.0,4.3) {total};
\draw[dashed, draw=engred!70] (4,0) -- (4,2.45);
\fill[engred!85] (4,2.45) circle (2.0pt);
\node[font=\scriptsize, anchor=south, text=engred!85] at (4,2.5) {$c^\star$};
\node[font=\small, anchor=south west] at (0.1, 3.7)
  {\textbf{(b) the staffing optimum}};
\end{scope}
\end{tikzpicture}
\caption[Queue model: validation against simulation and the staffing
optimum]{(a) The Erlang-C analytic prediction of mean wait against the
number of servers (solid line) compared with a discrete-event Monte-Carlo
simulation of the same arrival and service process (points, with one
standard error from $40$ replicated runs). The analytic curve falls
within the simulation's sampling interval at every staffing level, which
is the validation: two independent computations of the same quantity, one
closed-form and one stochastic, agree to within the simulation's noise.
(b) The optimisation the model exists to support. Staffing cost rises
linearly with $c$; expected waiting cost (lost time priced per customer
per minute) falls convexly as congestion eases. Their sum is U-shaped,
and its minimum $c^\star$ is the staffing level the model recommends. The
optimum sits at the staffing level where the marginal labour cost of one
more server equals the marginal waiting cost it removes, the discrete
analogue of the first-order condition Volume II Chapter 11 derives for
smooth objectives.}
\label{fig:vol02:ch18:queue-validation}
\end{figure}
